\begin{Answer}

\setcounter{ANSchapter}{1}
\answer
\begin{enumerate}
\item
以下に，章末問題解答に図表，数式が入る例を示します。
\begin{eqnarray}
 I_x & \= & I_0e^{-x/\delta} \notag \\
     & \= & a + b + c + d + e + f + g \notag \\
     & \leq & a + b + c + d + e + f + g \\
     & \geq & \Frac{1+2+3+4+5+6}{\sin\frac{\theta}{\alpha +\beta}}
\label{eqS.1}
\end{eqnarray}

\noindent
このように，章末問題解答で数式を書いた場合，式番は式\ref{eqS.1}のようになります。
図表を挿入した場合，図番は以下のように\FIGref{figS.4}となります。
表についても同様に，表番は\TABref{tabS.1}のように表記されます。

\begin{figure}[h]
\hfil
\includegraphics{square.eps}
\caption{これは章末問題解答内の図面です}
\label{figS.4}
\end{figure}

\begin{table}[h]\let\footnote\footnote
\begin{center}
\caption{表の例}
\label{tabS.1}
\begin{tabular}{l|l|r}
\Hline
項目1 & 項目2 & 項目3\\
\hline
a & A & 1\\
ab & AB & 12\\
abc & ABC & 123\\
\hline
\end{tabular}
\end{center}
\end{table}

\noindent
ここで，本文で式番を引用する場合は「式」の文字を手入力する必要がありますが，図表については\verb|\ref|で解図，解表の文字も自動的に表記されます。
なお，式番を本文中で引用する場合，\ref{eqS.1}式のようにはせず，必ず式\ref{eqS.1}とします。

\newpage
\item
これは，入れ子になった問題解答の例です。
\begin{enumerate}
\item
章末問題を作成する環境はenumerate環境を利用してい
ますので，各問題は``\verb@\item@''で始めて下さい。

\item
``\verb@\label@''および``\verb@\ref@''コマンドを使えば
\ref{CORONAex1}のように括弧付きで引用できます。
\end{enumerate}
\end{enumerate}


\end{Answer}
